% MT061 / explanation — canonical editable content.
题目问\textbf{预测函数占多少内存}，不是训练时需要保存多大的核矩阵。先把最终预测式写出来，再数需要存的数。


线性模型为\(f(x)=w^Tx+b\)。1000维输入需要1000个权重及1个偏置，共1001个数。即使训练中用了很多支持向量，它们也能合并为同一个\(w\)；所以\(C=0.1\)与\(C=1000\)在线性压缩表示下并列。


一般RBF模型为\(f(x)=\sum_{i=1}^{M}\beta_i k(x_i,x)+b\)，其中\(\beta_i=\alpha_i y_i\)，\(M\)是实际支持向量数。预测时要保存\(M\)个1000维向量、\(M\)个系数及1个偏置，即



\[
1000M+M+1=1001M+1\text{ 个数}.
\]


补充例子：若\(M=100\)，需要100,101个数；若\(M=500\)，需要500,501个数。因此常规非退化RBF模型比压缩后的线性模型存得多，两个RBF之间则要看谁的\(M\)小。数据类型相同且忽略软件开销时，数字个数才能直接用于内存比较。


原题只给\(C\)，没有给两个\(M\)，所以不足以普遍确定两种RBF的先后。配套答案假设大\(C\)的支持向量更少，但数量并不对\(C\)普遍单调。可靠作答是“两线性并列；RBF按实际支持向量数排序”，并注明原解使用的额外假设。


\textbf{出处：} 2023B期中 Q10，原卷第7页（10分）


\textbf{答案依据：} 2023B期中 Q10 随卷／配套参考答案；条件与解释经复核。
